\documentclass[a4paper]{article}
\usepackage[a4paper,margin=0mm]{geometry}
\usepackage{fontspec}
\usepackage{unicode-math}
\usepackage{tikz}
\usetikzlibrary{fit,calc,positioning,arrows.meta}
\usepackage[table]{xcolor}
\usepackage[absolute]{textpos}
\usepackage{array}

\setmainfont{Libertinus Serif}
\setsansfont{Libertinus Sans}
\setmathfont{Libertinus Math}

\definecolor{accent}{HTML}{0F6A52}
\definecolor{badge}{HTML}{A9CFC2}
\definecolor{ink}{HTML}{2A2E30}

\setlength{\TPHorizModule}{1mm}
\setlength{\TPVertModule}{1mm}
\setlength{\parindent}{0pt}
\pagestyle{empty}

\newcommand{\topicheader}[2]{%
\begin{textblock*}{194mm}(16mm,10mm)
\begin{tikzpicture}
\node[fill=accent,text=white,font=\sffamily\bfseries\fontsize{13}{13}\selectfont,
      minimum width=9mm,minimum height=9mm,inner sep=0pt] at (4.5mm,-4.5mm) {#1};
\node[font=\sffamily\bfseries\fontsize{15}{16}\selectfont,text=ink,anchor=west] at (12mm,-4.5mm) {#2};
\end{tikzpicture}
\end{textblock*}
\begin{textblock*}{194mm}(16mm,22.5mm)
\begin{tikzpicture}
\fill[accent] (0,0) rectangle (194mm,0.5mm);
\end{tikzpicture}
\end{textblock*}
}

\newcommand{\pagekicker}[1]{%
\begin{textblock*}{100mm}(16mm,25.5mm)
\begin{tikzpicture}
\node[fill=ink,text=white,font=\sffamily\bfseries\fontsize{8}{8}\selectfont,
      inner sep=2.2pt,rounded corners=1pt] {\MakeUppercase{#1}};
\end{tikzpicture}
\end{textblock*}
}

\newcommand{\videobanner}[1]{%
\begin{textblock*}{178mm}(16mm,25.5mm)
\begin{tikzpicture}
\fill[accent] (0,0) rectangle (178mm,-10mm);
\draw[white,line width=0.6pt] (3mm,-1.7mm) rectangle (11.5mm,-8.3mm);
\node[white,font=\sffamily\fontsize{4}{4.4}\selectfont,align=center] at (7.25mm,-5mm) {QR\\CODE};
\node[white,font=\sffamily\bfseries\fontsize{9.5}{11}\selectfont,anchor=west] at (16mm,-5.6mm) {WATCH THE VIDEO BEFORE YOU SOLVE};
\node[white,font=\sffamily\fontsize{7.5}{8}\selectfont,anchor=east] at (174mm,-5.6mm) {YouTube ~\textbullet~ #1};
\end{tikzpicture}
\end{textblock*}
}

\newcommand{\pagefooter}[1]{%
\begin{textblock*}{194mm}(8mm,286.41mm)
  \begin{tikzpicture}
    \fill[accent] (0,0) rectangle (194mm,0.6mm);
  \end{tikzpicture}
\end{textblock*}
\begin{textblock*}{63.59mm}(8mm,288.27mm)
  {\sffamily\bfseries\color{accent}\fontsize{13}{16}\selectfont SAT 1200+}
\end{textblock*}
\begin{textblock*}{10mm}(188mm,288.27mm)
  \centering{\sffamily\color{accent}\fontsize{14}{16}\selectfont #1}
\end{textblock*}
}

\newcommand{\cardbg}[4]{%
\begin{textblock*}{#3mm}(#1mm,#2mm)
\begin{tikzpicture}
\fill[badge!22,rounded corners=3mm] (0,0) rectangle (#3mm,-#4mm);
\end{tikzpicture}
\end{textblock*}%
}

\newcommand{\pill}[3]{%
\begin{textblock*}{160mm}(#1mm,#2mm)
\begin{tikzpicture}
\node[fill=accent,text=white,anchor=north west,font=\sffamily\bfseries\fontsize{10}{10}\selectfont,
      rounded corners=8pt,inner xsep=5mm,inner ysep=1.8mm] at (0,0) {RULE ~\textperiodcentered~ #3};
\end{tikzpicture}
\end{textblock*}%
}

\newcommand{\pillw}[3]{%
\begin{textblock*}{160mm}(#1mm,#2mm)
\begin{tikzpicture}
\node[fill=accent,text=white,anchor=north west,font=\sffamily\bfseries\fontsize{10}{10}\selectfont,
      rounded corners=8pt,inner xsep=5mm,inner ysep=1.8mm] at (0,0) {#3};
\end{tikzpicture}
\end{textblock*}%
}

\arrayrulecolor{accent}
\setlength{\arrayrulewidth}{0.7pt}
\renewcommand{\arraystretch}{1.5}


% Bar graph (histogram of counts). #1 counts, #2 slot width mm, #3 mm per count,
% #4 number of bars, #5 total picture height mm, #6 mean (x position), #7 top of mean line (count units)
\newcommand{\bargraph}[7]{%
\pgfmathsetmacro{\bgW}{#4*#2}%
\pgfmathsetmacro{\bgS}{#5-5}%
\begin{tikzpicture}[x=#2mm,y=#3mm]
\useasboundingbox (0mm,-#5mm) rectangle (\bgW mm,0mm);
\begin{scope}[shift={(0mm,-\bgS mm)}]
\draw[ink!55,line width=0.5pt] (0.3,0) -- ({#4+0.7},0);
\foreach \c [count=\i] in {#1} {%
  \ifnum\c>0 \fill[accent] ({\i-0.42},0) rectangle ({\i+0.42},\c);\fi
  \node[below,font=\sffamily\fontsize{6}{6}\selectfont,text=ink!80,inner sep=1pt] at (\i,0) {\i};
}
\draw[dashed,ink!70,line width=0.45pt] (#6,-0.2) -- (#6,#7);
\node[above,font=\sffamily\fontsize{5.5}{6}\selectfont,text=ink!80,inner sep=1pt] at (#6,#7) {mean};
\end{scope}
\end{tikzpicture}}

\newwrite\mw \immediate\openout\mw=meas.txt
\begin{document}
\setbox0=\vbox{\hsize=158mm\noindent {\fontsize{15}{18}\selectfont mean $=\dfrac{\text{total}}{\text{count}}$ \quad median $=$ middle value \quad mode $=$ most frequent}\par}\immediate\write\mw{0 \the\dimexpr\ht0+\dp0\relax}
\setbox0=\vbox{\hsize=158mm\noindent {\fontsize{13}{16}\selectfont even count: median $=$ average of the two middle values}\par}\immediate\write\mw{1 \the\dimexpr\ht0+\dp0\relax}
\setbox0=\vbox{\hsize=158mm\noindent {\fontsize{13}{16}\selectfont mean of 5 is $30$, mean of the other 4 is $28$: \ removed $=5(30)-4(28)=38$}\par}\immediate\write\mw{2 \the\dimexpr\ht0+\dp0\relax}
\setbox0=\vbox{\hsize=158mm\noindent {\itshape\color{black!70}\fontsize{11}{14}\selectfont Work with totals, not means.}\par}\immediate\write\mw{3 \the\dimexpr\ht0+\dp0\relax}
\setbox0=\vbox{\hsize=158mm\noindent {\fontsize{15}{18}\selectfont extreme value: mean and range move; median barely; mode not}\par}\immediate\write\mw{4 \the\dimexpr\ht0+\dp0\relax}
\setbox0=\vbox{\hsize=158mm\noindent {\fontsize{13}{16}\selectfont $3,5,5,8$ plus $30$: \ mean $5.25\to10.2$, \ median $5\to5$, \ mode $5\to5$, \ range $5\to27$}\par}\immediate\write\mw{5 \the\dimexpr\ht0+\dp0\relax}
\setbox0=\vbox{\hsize=158mm\noindent {\itshape\color{black!70}\fontsize{11}{14}\selectfont An even count makes the median the average of two middle values, so it can move. Compare totals before and after.}\par}\immediate\write\mw{6 \the\dimexpr\ht0+\dp0\relax}
\setbox0=\vbox{\hsize=158mm\noindent {\fontsize{15}{18}\selectfont mean $=\dfrac{\sum(\text{value}\times\text{frequency})}{\text{total frequency}}$ \qquad combined: $\dfrac{n_1m_1+n_2m_2}{n_1+n_2}$}\par}\immediate\write\mw{7 \the\dimexpr\ht0+\dp0\relax}
\setbox0=\vbox{\hsize=158mm\noindent {\fontsize{13}{16}\selectfont 15 students (mean 60) and 25 students (mean 80): \ $\dfrac{900+2000}{40}=72.5$}\par}\immediate\write\mw{8 \the\dimexpr\ht0+\dp0\relax}
\setbox0=\vbox{\hsize=158mm\noindent {\itshape\color{black!70}\fontsize{11}{14}\selectfont Missing value: needed $=$ (target mean $\times$ new count) $-$ (current total). Never average two means without weights.}\par}\immediate\write\mw{9 \the\dimexpr\ht0+\dp0\relax}
\setbox0=\vbox{\hsize=158mm\noindent {\fontsize{13}{16}\selectfont SD $=$ how far values typically sit from the mean. \textbf{More spread $\Rightarrow$ bigger SD.}}\par}\immediate\write\mw{10 \the\dimexpr\ht0+\dp0\relax}
\setbox0=\vbox{\hsize=158mm\noindent {\fontsize{13}{16}\selectfont skyrocket bar near the mean $\Rightarrow$ \textbf{small SD} \quad double top at the ends $\Rightarrow$ \textbf{large SD}}\par}\immediate\write\mw{11 \the\dimexpr\ht0+\dp0\relax}
\setbox0=\vbox{\hsize=158mm\noindent {\itshape\color{black!70}\fontsize{11}{14}\selectfont All three graphs have the same mean. Judge spread by where the tall bars sit, not by the range alone.}\par}\immediate\write\mw{12 \the\dimexpr\ht0+\dp0\relax}
\setbox0=\vbox{\hsize=158mm\noindent {\fontsize{10.5}{13}\selectfont\renewcommand{\arraystretch}{1.25}
\begin{tabular}{|p{57mm}|p{25mm}|p{26mm}|p{30mm}|}\hline
\cellcolor{accent}\textcolor{white}{\bfseries Change to the data} & \cellcolor{accent}\textcolor{white}{\bfseries Mean} & \cellcolor{accent}\textcolor{white}{\bfseries SD} & \cellcolor{accent}\textcolor{white}{\bfseries Range, IQR}\\ \hline
add $c$ to every value & $+c$ & same & same\\ \hline
multiply every value by $k$ & $\times k$ & $\times|k|$ & $\times|k|$\\ \hline
add a far-away outlier & pulled & jumps up & range jumps\\ \hline
add a value equal to the mean & same & smaller & same\\ \hline
\end{tabular}}\par}\immediate\write\mw{13 \the\dimexpr\ht0+\dp0\relax}
\setbox0=\vbox{\hsize=158mm\noindent {\fontsize{15}{18}\selectfont range $=\max-\min$ \qquad IQR $=Q_3-Q_1$ (the middle half)}\par}\immediate\write\mw{14 \the\dimexpr\ht0+\dp0\relax}
\setbox0=\vbox{\hsize=158mm\noindent {\fontsize{13}{16}\selectfont SD $=$ typical distance from the mean \qquad median and IQR resist outliers}\par}\immediate\write\mw{15 \the\dimexpr\ht0+\dp0\relax}
\setbox0=\vbox{\hsize=158mm\noindent {\fontsize{13}{16}\selectfont two clumps far apart: the combined data have a larger SD than either clump alone}\par}\immediate\write\mw{16 \the\dimexpr\ht0+\dp0\relax}
\setbox0=\vbox{\hsize=158mm\noindent {\itshape\color{black!70}\fontsize{11}{14}\selectfont Split into halves to find $Q_1$ and $Q_3$ (medians of the lower and upper half). Range and IQR can disagree.}\par}\immediate\write\mw{17 \the\dimexpr\ht0+\dp0\relax}
\setbox0=\vbox{\hsize=158mm\noindent {\fontsize{13}{16}\selectfont box plot = five numbers: each quarter of the data lies between neighbors. \textbf{No mean.}}\par}\immediate\write\mw{18 \the\dimexpr\ht0+\dp0\relax}
\setbox0=\vbox{\hsize=158mm\noindent {\fontsize{13}{16}\selectfont histogram: height $=$ frequency; median bar $=$ where the running total hits half}\par}\immediate\write\mw{19 \the\dimexpr\ht0+\dp0\relax}
\setbox0=\vbox{\hsize=158mm\noindent {\itshape\color{black!70}\fontsize{11}{14}\selectfont A long tail on the right: mode $<$ median $<$ mean. Symmetric: mean $=$ median.}\par}\immediate\write\mw{20 \the\dimexpr\ht0+\dp0\relax}
\setbox0=\vbox{\hsize=158mm\noindent {\fontsize{15}{18}\selectfont least mean: every value at the \textbf{bottom} of its class; greatest: at the \textbf{top}}\par}\immediate\write\mw{21 \the\dimexpr\ht0+\dp0\relax}
\setbox0=\vbox{\hsize=158mm\noindent {\fontsize{13}{16}\selectfont integer classes $0$--$9$ and $10$--$19$ (5 values each): least mean $5$, greatest $14$, difference $9$}\par}\immediate\write\mw{22 \the\dimexpr\ht0+\dp0\relax}
\setbox0=\vbox{\hsize=158mm\noindent {\itshape\color{black!70}\fontsize{11}{14}\selectfont For the median, find the two middle positions first (even count), then bound each. Midpoints give an estimate, not a bound.}\par}\immediate\write\mw{23 \the\dimexpr\ht0+\dp0\relax}
\setbox0=\vbox{\hsize=158mm\noindent {\fontsize{15}{18}\selectfont $y=mx+b$: \ slope $=$ change in $y$ per \textbf{one} unit of $x$}\par}\immediate\write\mw{24 \the\dimexpr\ht0+\dp0\relax}
\setbox0=\vbox{\hsize=158mm\noindent {\fontsize{13}{16}\selectfont $b=$ predicted $y$ at $x=0$ \qquad residual $=$ actual $-$ predicted}\par}\immediate\write\mw{25 \the\dimexpr\ht0+\dp0\relax}
\setbox0=\vbox{\hsize=158mm\noindent {\fontsize{13}{16}\selectfont $y=1.5x+8$, point $(12,30)$: \ residual $30-26=+4$ \quad ($+$ means above the line)}\par}\immediate\write\mw{26 \the\dimexpr\ht0+\dp0\relax}
\setbox0=\vbox{\hsize=158mm\noindent {\itshape\color{black!70}\fontsize{11}{14}\selectfont Inside the data range: prediction. Far outside it: extrapolation, unreliable. Check the intercept at $x=0$, not the left edge.}\par}\immediate\write\mw{27 \the\dimexpr\ht0+\dp0\relax}
\setbox0=\vbox{\hsize=158mm\noindent {\fontsize{15}{18}\selectfont slope from two far-apart points: $\dfrac{y_2-y_1}{x_2-x_1}$ \qquad point $(x_0,y_0)$: $b=y_0-mx_0$}\par}\immediate\write\mw{28 \the\dimexpr\ht0+\dp0\relax}
\setbox0=\vbox{\hsize=158mm\noindent {\fontsize{13}{16}\selectfont exponential $y=ab^x$: equal steps in $x$ multiply $y$ by $b$ \quad $4,\,6,\,9,\,13.5\Rightarrow b=1.5$, $a=4$}\par}\immediate\write\mw{29 \the\dimexpr\ht0+\dp0\relax}
\setbox0=\vbox{\hsize=158mm\noindent {\itshape\color{black!70}\fontsize{11}{14}\selectfont A far-away point tilts the line toward itself; removing it moves the slope and the intercept in opposite ways.}\par}\immediate\write\mw{30 \the\dimexpr\ht0+\dp0\relax}
\setbox0=\vbox{\hsize=158mm\noindent {\fontsize{15}{18}\selectfont $P(A\mid B)=\dfrac{\text{count}(A\text{ and }B)}{\text{count}(B)}$ \quad the ``given'' group is the denominator}\par}\immediate\write\mw{31 \the\dimexpr\ht0+\dp0\relax}
\setbox0=\vbox{\hsize=158mm\noindent {\fontsize{13}{16}\selectfont 40 like tea, 10 of them under 30: \ $P(\text{under }30\mid\text{tea})=\dfrac{10}{40}$ \qquad $P(\text{tea}\mid\text{under }30)\ne$ that}\par}\immediate\write\mw{32 \the\dimexpr\ht0+\dp0\relax}
\setbox0=\vbox{\hsize=158mm\noindent {\itshape\color{black!70}\fontsize{11}{14}\selectfont Reverse conditionals: with percents, build both parts (\,$0.55\cdot0.30$ and $0.45\cdot0.60$\,) and divide by the total of the group asked about.}\par}\immediate\write\mw{33 \the\dimexpr\ht0+\dp0\relax}
\setbox0=\vbox{\hsize=158mm\noindent {\fontsize{15}{18}\selectfont independent $\iff P(A\mid B)=P(A)\iff P(A\text{ and }B)=P(A)\,P(B)$}\par}\immediate\write\mw{34 \the\dimexpr\ht0+\dp0\relax}
\setbox0=\vbox{\hsize=158mm\noindent {\fontsize{13}{16}\selectfont in a table: the rate is the same in every row, e.g. $\dfrac{12}{40}=\dfrac{60}{200}=0.3$ \ $\Rightarrow$ independent}\par}\immediate\write\mw{35 \the\dimexpr\ht0+\dp0\relax}
\setbox0=\vbox{\hsize=158mm\noindent {\itshape\color{black!70}\fontsize{11}{14}\selectfont A parameter in the table: set the two rates equal and solve.}\par}\immediate\write\mw{36 \the\dimexpr\ht0+\dp0\relax}
\setbox0=\vbox{\hsize=158mm\noindent {\fontsize{15}{18}\selectfont $P(\text{at least one})=1-P(\text{none})$ \qquad independent: multiply}\par}\immediate\write\mw{37 \the\dimexpr\ht0+\dp0\relax}
\setbox0=\vbox{\hsize=158mm\noindent {\fontsize{13}{16}\selectfont $4$ items, each defective with probability $0.05$: \ $1-0.95^4$}\par}\immediate\write\mw{38 \the\dimexpr\ht0+\dp0\relax}
\setbox0=\vbox{\hsize=158mm\noindent {\fontsize{13}{16}\selectfont without replacement the pool shrinks: \ $5$ red, $3$ blue: $P(\text{both red})=\dfrac58\cdot\dfrac47$}\par}\immediate\write\mw{39 \the\dimexpr\ht0+\dp0\relax}
\setbox0=\vbox{\hsize=158mm\noindent {\itshape\color{black!70}\fontsize{11}{14}\selectfont Add the disjoint cases (e.g. both red, both blue). For ``least $n$ so that\,\ldots'' try values.}\par}\immediate\write\mw{40 \the\dimexpr\ht0+\dp0\relax}
\setbox0=\vbox{\hsize=158mm\noindent {\fontsize{15}{18}\selectfont expected count $=$ proportion $\times$ population}\par}\immediate\write\mw{41 \the\dimexpr\ht0+\dp0\relax}
\setbox0=\vbox{\hsize=158mm\noindent {\fontsize{13}{16}\selectfont 36 of 200 sampled $=18\%$; town of $5000$: \ $0.18\cdot5000=900$}\par}\immediate\write\mw{42 \the\dimexpr\ht0+\dp0\relax}
\setbox0=\vbox{\hsize=158mm\noindent {\itshape\color{black!70}\fontsize{11}{14}\selectfont Chain the proportions by multiplying. Percent of a subgroup: part $\div$ that subgroup's total, not the grand total.}\par}\immediate\write\mw{43 \the\dimexpr\ht0+\dp0\relax}
\setbox0=\vbox{\hsize=158mm\noindent {\fontsize{15}{18}\selectfont estimate $\pm$ margin of error $=$ \textbf{plausible} values for the population}\par}\immediate\write\mw{44 \the\dimexpr\ht0+\dp0\relax}
\setbox0=\vbox{\hsize=158mm\noindent {\fontsize{13}{16}\selectfont $62\%$ with margin of error $3$ points: \ plausible $59\%$ to $65\%$ (not exactly $62\%$)}\par}\immediate\write\mw{45 \the\dimexpr\ht0+\dp0\relax}
\setbox0=\vbox{\hsize=158mm\noindent {\itshape\color{black!70}\fontsize{11}{14}\selectfont 4 times the sample size halves the margin. Overlapping intervals: no difference shown. It is not a chance of error.}\par}\immediate\write\mw{46 \the\dimexpr\ht0+\dp0\relax}
\setbox0=\vbox{\hsize=158mm\noindent {\fontsize{10.5}{13}\selectfont\renewcommand{\arraystretch}{1.25}
\begin{tabular}{|p{66mm}|p{74mm}|}\hline
\cellcolor{accent}\textcolor{white}{\bfseries Study design} & \cellcolor{accent}\textcolor{white}{\bfseries Conclusion allowed}\\ \hline
random sample \textbf{and} random assignment & cause for the whole population\\ \hline
random sample, no random assignment & association for the population (no cause)\\ \hline
random assignment of volunteers & cause, for people like the volunteers\\ \hline
neither (self-selected, convenience) & association in the sample; biased\\ \hline
\end{tabular}}\par}\immediate\write\mw{47 \the\dimexpr\ht0+\dp0\relax}
\setbox0=\vbox{\hsize=158mm\noindent {\itshape\color{black!70}\fontsize{11}{14}\selectfont Generalize only to the population the sample was drawn from. A flawed sampling method cannot be fixed by a larger sample.}\par}\immediate\write\mw{48 \the\dimexpr\ht0+\dp0\relax}
\immediate\closeout\mw
\end{document}